\chapter{Organizzazione del codice}
\label{cap:codebase}

Questo capitolo descrive la struttura del codice e i meccanismi che ne
garantiscono la coerenza. Serve come mappa per chi voglia riprodurre o estendere
il lavoro, e documenta le scelte organizzative che hanno effettivamente
prevenuto guasti.

\section{Struttura dei moduli}

\begin{center}
\begin{tabular}{ll}
\toprule
Percorso & Responsabilità \\
\midrule
\texttt{src/datasets/} & caricamento, suddivisione, retrieval, transizioni \\
\texttt{src/models/} & caricamento del modello, teste aggiuntive \\
\texttt{src/grammar/} & Trie, maschera, processore di logits \\
\texttt{src/rewards/t2g\_rewards.py} & le componenti di reward \\
\texttt{src/training/} & punti di ingresso e trainer (SFT, GRPO, ausiliari) \\
\texttt{src/analysis/} & baseline a regole, diagnostiche \\
\texttt{src/utils/} & metriche, configurazione, monitoraggio \\
\texttt{cluster/} & script SLURM e libreria di shell \\
\texttt{remote/} & servizio FastAPI e interfaccia testuale \\
\texttt{tests/} & suite di test ($409$ test) \\
\texttt{experiments/configs/qwen25-05b/} & configurazioni delle celle \\
\bottomrule
\end{tabular}
\end{center}

Due note sulla separazione delle responsabilità.

La prima: \texttt{src/grammar/} è indipendente dal modello. Il Trie e la
maschera sono strutture dati pure, e tutti i test che ne verificano la
correttezza girano senza GPU. Questo ha reso possibile sviluppare e verificare
il vincolo di decodifica sulla macchina locale, dove l'addestramento non è
eseguibile.

La seconda: \texttt{src/analysis/} contiene codice che non partecipa
all'addestramento ma è essenziale all'interpretazione. La baseline a regole vive
qui, non fra i modelli, perché non è un modello: è uno strumento di misura.

\section{Il sistema di configurazione}

Ogni cella sperimentale è descritta da un file YAML. Il meccanismo di
composizione è la chiave \texttt{extends}, risolta da
\texttt{src/utils/config.py}, funzione \texttt{resolve\_config}, che esegue una
\emph{fusione profonda ricorsiva}: i dizionari annidati vengono fusi chiave per
chiave, invece di essere sostituiti in blocco.

La distinzione importa. Con una sostituzione in blocco, una cella che volesse
cambiare un solo iperparametro dell'ottimizzatore dovrebbe ridichiarare tutta la
sezione, e ogni modifica al file di base non si propagherebbe. Con la fusione
profonda, una cella dichiara solo le differenze.

Un dettaglio di igiene: la chiave \texttt{extends} viene rimossa dalla
configurazione risolta e non raggiunge i trainer. Questi ricevono una
configurazione già completa e piatta, e non hanno modo di sapere da quale
gerarchia provenga. È una separazione voluta: la logica di composizione sta in un
solo posto.

\section{La matrice sperimentale}

Il progetto definisce \textbf{15 celle}, articolate in \textbf{27 voci} di
esecuzione (una cella di addestramento genera tipicamente una voce di
addestramento e una di valutazione).

\begin{center}
\begin{tabular}{lll}
\toprule
Gruppo & Cella & Tipo \\
\midrule
\texttt{baseline} & \texttt{zero-shot} & sola valutazione \\
                  & \texttt{zero-shot-no-grammar} & sola valutazione \\
                  & \texttt{few-shot} & sola valutazione \\
\midrule
\texttt{sft} & \texttt{zero-shot} & addestramento \\
\midrule
\texttt{grpo} & \texttt{zero-shot} & addestramento \\
              & \texttt{few-shot} & addestramento \\
\midrule
\texttt{sft-grpo} & \texttt{zero-shot} & addestramento \\
                  & \texttt{few-shot} & addestramento \\
\midrule
\texttt{ablations/rewards} & \texttt{edit-validity} & addestramento \\
                           & \texttt{historical-stack} & addestramento \\
\texttt{ablations/loss} & \texttt{dr-grpo} & addestramento \\
\texttt{ablations/decoding} & \texttt{no-grammar} & addestramento \\
                            & \texttt{hot-rollout} & addestramento \\
\texttt{ablations/objectives} & \texttt{sft-allowed-mass} & addestramento \\
                              & \texttt{sft-structured} & addestramento \\
\bottomrule
\end{tabular}
\end{center}

Le tre celle \texttt{baseline} sono di sola valutazione: misurano il modello non
addestrato. Le loro configurazioni non contengono la sezione di addestramento, ed
è precisamente questa asimmetria che ha prodotto il guasto documentato nella
sezione~\ref{sec:bug-evalonly}.

Le cinque celle di ablazione sono le più interessanti dal punto di vista
metodologico, perché ciascuna isola un fattore che la parte principale
dell'esperimento confonde:
\begin{itemize}
  \item \texttt{edit-validity} contro \texttt{historical-stack} isola l'effetto
        della revisione dello stack di reward (Capitolo~\ref{cap:reward});
  \item \texttt{dr-grpo} attiva esplicitamente la variante di denominatore che
        nessuna configurazione storica selezionava
        (sezione~\ref{sec:denominatori});
  \item \texttt{no-grammar} isola il contributo del vincolo simbolico
        (Capitolo~\ref{cap:decoding});
  \item \texttt{hot-rollout} varia la temperatura di generazione dei candidati,
        cioè il parametro che governa direttamente il supporto dei rollout
        (sezione~\ref{sec:supporto});
  \item le due celle in \texttt{objectives} corrispondono ai due obiettivi
        ausiliari del Capitolo~\ref{cap:ausiliari}.
\end{itemize}

\section{La validazione delle configurazioni}

Il file \texttt{tests/validate\_configs.py} è il punto in cui le invarianti
scoperte durante il progetto sono codificate. Non è un test di stile: ogni
controllo corrisponde a un guasto reale o a un attacco dimostrato.

\begin{center}
\begin{tabular}{ll}
\toprule
Controllo & Motivazione \\
\midrule
presenza delle sezioni attese & un errore di struttura deve fallire subito \\
tipi dei valori & impedisce coercizioni silenziose \\
somma dei pesi di reward $= 1{,}0$ & la scala della reward non deve variare fra celle \\
chiavi morte rifiutate & impedisce configurazioni che sembrano attive e non lo sono \\
knob dell'obiettivo RL validi & rende esplicita la variante ottimizzata \\
peso del termine fuori vocabolario in $[0;\,0{,}6]$ & sopra $0{,}6$ esiste un attacco dimostrato \\
\texttt{max\_prompt\_length} $\ge 512$ con retrieval & un prompt troncato annulla il few-shot \\
\bottomrule
\end{tabular}
\end{center}

Vale la pena soffermarsi su due voci.

\subsection{Le chiavi morte}

Durante la pulizia del progetto sono state rimosse tre componenti di reward
(sezione~\ref{sec:reward-rimosse}) e diversi percorsi legacy. Il rischio, dopo
una rimozione, è che una configurazione continui a impostare una chiave che non
esiste più: il file appare configurato in un certo modo, il codice ignora la
chiave, e l'esperimento gira con impostazioni diverse da quelle dichiarate.

Il validatore rifiuta esplicitamente le chiavi rimosse. È una forma di
protezione contro la classe di guasto più costosa in un contesto sperimentale:
la discrepanza fra configurazione dichiarata e comportamento effettivo.

\subsection{Il vincolo sulla lunghezza del prompt}

Il controllo \texttt{max\_prompt\_length} $\ge 512$ in presenza di retrieval
merita una nota, perché protegge da un fallimento perfettamente silenzioso. Se
il prompt viene troncato, gli esempi recuperati vengono eliminati, e la cella
few-shot diventa indistinguibile da una cella zero-shot. Nessun errore, nessun
avviso: solo numeri che sembrano dire che il few-shot non serve.

Dato che il fattore prompting è, secondo la
sezione~\ref{sec:decomposizione}, l'effetto più grande misurato in tutto il
progetto, un guasto che lo annullasse silenziosamente sarebbe particolarmente
grave.

\section{Le diagnostiche di Markov}

Il modulo \texttt{src/analysis/markov\_diagnostics.py} implementa strumenti di
ispezione sui percorsi di transizione fra gloss:

\begin{center}
\begin{tabular}{ll}
\toprule
Funzione & Ruolo \\
\midrule
\texttt{path\_log\_energy} & punteggio logaritmico di un percorso \\
\texttt{hard\_viterbi\_diagnostic} & percorso di massima verosimiglianza \\
\texttt{soft\_viterbi\_diagnostic} & versione smussata (somma-log-esponenziale) \\
\bottomrule
\end{tabular}
\end{center}

Due scelte implementative vanno segnalate.

La prima: le funzioni sono validate contro un calcolo a forza bruta su istanze
piccole. Un algoritmo di programmazione dinamica come Viterbi è facile da
implementare in modo sottilmente sbagliato (un indice fuori posto produce
risultati plausibili), e il confronto con l'enumerazione esaustiva su casi di
dimensione ridotta è la verifica che coglie quel tipo di errore.

La seconda: esiste una guardia che rifiuta spazi di stato superiori a $256$
elementi. È un limite deliberato, non un difetto: queste funzioni servono a
ispezionare casi specifici, e permetterne l'uso su spazi grandi le trasformerebbe
in un collo di bottiglia computazionale mascherato da strumento di analisi.

Il punto più importante è lo \emph{statuto} di questo modulo. Le quantità di
Viterbi sono state \textbf{rimosse dalle reward}
(sezione~\ref{sec:reward-rimosse}) perché a lunghezza uguale collassano su un
segnale già presente. Sono invece state \textbf{conservate come diagnostica}. La
distinzione è sostanziale: una quantità può essere informativa da ispezionare e
inutile, o addirittura dannosa, da ottimizzare.

\section{Il servizio remoto}

La directory \texttt{remote/} contiene un servizio FastAPI e un'interfaccia
testuale per interagire con il cluster: sottomettere celle, ispezionare la coda,
recuperare risultati.

La sua ragione d'essere è il vincolo descritto nel
Capitolo~\ref{cap:infrastruttura}: la macchina di sviluppo non può eseguire
l'addestramento, e il nodo di login non ha il runtime del progetto. Ogni
operazione sperimentale passa dunque per una connessione remota, e avere
un'interfaccia unica riduce la superficie di errore rispetto a comandi
costruiti a mano.

\section{La suite di test}

Lo stato corrente è di $409$ test. La loro distribuzione riflette la strategia
adottata: la maggior parte verifica componenti eseguibili senza GPU.

\begin{description}
  \item[Strutture simboliche.] Il Trie, la maschera, il calcolo dell'insieme
        ammesso. Nessuna dipendenza dal modello.
  \item[Reward.] Ogni componente valutata su candidati costruiti a mano,
        compresi quelli avversari del Capitolo~\ref{cap:reward}.
  \item[Metriche.] Compresi i casi che documentano i difetti di ROUGE-L
        (sezione~\ref{sec:difetto-rouge}): quei valori sono asseriti nei test, in
        modo che il difetto sia registrato e non possa essere ``corretto'' per
        distrazione rendendo i numeri non confrontabili.
  \item[Contratti degli obiettivi ausiliari.] L'inerzia esatta a peso nullo e i
        rifiuti espliciti delle configurazioni incompatibili
        (sezione~\ref{sec:contratto-ausiliari}).
  \item[Guardie dei punti di ingresso.] I dieci test della
        sezione~\ref{sec:bug-evalonly}, con l'invariante diretta e quella
        inversa.
\end{description}

\subsection{Un test rosso dichiarato}

Un test risulta fallito nella suite corrente:
\texttt{tests/test\_remote\_app.py::test\_ssh\_alias\_no\_user\_no\_port}.

Il fallimento è stato verificato presente anche sul codice non modificato:
non è una regressione introdotta dal lavoro descritto in questa relazione. La
causa è che il test effettua una connessione SSH reale verso il cluster, quindi
il suo esito dipende dalla disponibilità della rete e dalla configurazione
dell'ambiente di esecuzione, non dal codice.

È riportato qui, invece di essere silenziato, perché lo stato della suite fa
parte dei risultati: dichiarare ``409 test, tutti verdi'' quando uno è rosso
sarebbe una piccola falsificazione, e un test che dipende dalla rete è un difetto
di disegno del test da correggere separatamente, non da nascondere.

\section{La pulizia del codice}

Una parte sostanziale del lavoro è consistita nella rimozione di codice. Le voci
principali:

\begin{center}
\begin{tabular}{lr}
\toprule
Elemento rimosso & Note \\
\midrule
Automa a stack e grammatica LL(1) & mai eseguito, sovradimensionato (sez.~\ref{sec:pda-scartato}) \\
Tre componenti di reward basate su Viterbi & ridondanti (sez.~\ref{sec:reward-rimosse}) \\
Infrastruttura di supporto & circa $8\,774$ righe complessive \\
Percorsi di checkpoint legacy & puntavano a directory non più prodotte \\
Configurazioni obsolete & sostituite dalla matrice corrente \\
Sei documenti di progetto & superati o contraddittori \\
\bottomrule
\end{tabular}
\end{center}

La rimozione di codice funzionante richiede giustificazione, e in ciascun caso è
stata fornita: l'automa a stack non era mai stato eseguito e aveva un conflitto
di parsing irrisolto; le reward basate su Viterbi erano dimostrabilmente
ridondanti sia algebricamente sia empiricamente (entro $\pm 0{,}006$ dal
controllo); i percorsi legacy puntavano a directory che nessuna cella corrente
produce.

Il criterio applicato è che il codice non eseguito è un costo senza
contropartita: va mantenuto solo se esiste un piano concreto per eseguirlo. Gli
obiettivi ausiliari del Capitolo~\ref{cap:ausiliari} sono l'eccezione
deliberata, e per questo sono accompagnati da criteri di promozione espliciti.
